\documentclass[11pt,a4paper]{article}
\usepackage[margin=25mm]{geometry}
\usepackage{amsmath,amssymb,amsthm,booktabs}
\newtheorem{lemma}{Lemma}\newtheorem{theorem}{Theorem}
\title{Two order ten counts from an affine complement}
\author{Human Compression project}\date{}
\begin{document}\maketitle
Let $G=O'N$, with $|G|=460815505920$, and write
\[
 x_J(p,q)=|\{(r,s)\in J^2:|r|=p,\ |s|=q,\ |rs|=2\}|.
\]
Let $e_J(p,q)$ count those pairs that generate $J$. We prove the
fixed source entries $(5,10)$ and $(10,10)$. The complete subgroup
roster, primary ATLAS structures and acquired character tables are
borrowed premises. All row numbers refer to that roster.

We reuse two accepted facts. Every element of order ten lies in a
unique row 7 maximal subgroup $A$, an involution centralizer, and a
unique row 11 maximal subgroup $B$. A row 32 subgroup $H$ has order
2880 and has unique such maximal overgroups, whose intersection is
$H$. Moreover, $H$ has a normal subgroup $N\cong A_6$ with
$H/N\cong Q_8$; the preimage of the center of this quotient is
$A_6\times C_2$. The proofs of these facts are retained in the
accepted order 20 and earlier order 10 components.

\begin{lemma}\label{cover}
For either of our two types, a pair outside all subgroups in rows 7
and 11 belongs to exactly one maximal subgroup. It is either a
row 6 subgroup $J\cong J_1$ or a row 12 subgroup $M$ of order 25920.
\end{lemma}
\begin{proof}
Intersect the complete ancestor lists of the cyclic orders five
and ten, or of ten with itself, and remove the descendant lists of
rows 7 and 11. Apart from $G$ itself, the possible generated
subgroup rows in either case are exactly
\[
 \begin{array}{c|rrrrrr}
 \text{row}&6&12&20&52&90&286\\
 \text{order}&175560&25920&12960&1620&810&110.
 \end{array}
\]
Row 286 has just one immediate overgroup type, row 6. The direct
incidence $[G:J]\,1596=[G:C_{11}:C_{10}]$ shows that its actual
overgroup is unique. This also proves uniqueness for pairs
generating row 6.

The primary structure of row 12 is $M=V\rtimes K$, where
$V\cong C_3^4$ and $K$ has order 320, with a normal extraspecial
subgroup $E$ of order 32 and quotient $D_{10}$. In particular
$M/V$ has order prime to three. If a subgroup $L$ of one of the
orders 12960, 1620 or 810 lies in $M$, then its Sylow 3-subgroup
has order 81 and lies in $V$. Thus $V\leq L$, and $V$ is the unique
Sylow 3-subgroup of $L$. This subgroup is independent of the
chosen overgroup $M$. The sole row of order 81 is row 311, with
normalizer order 25920. Consequently every row 12 overgroup of
$L$ equals $N_G(V)$, and is unique. The complete ancestor lists
give no other maximal type for these three rows. This proves the
claim. Only single incidence counts and normalizer orders have
been used, not products of cover multiplicities.
\end{proof}

\begin{lemma}\label{H}
The numbers of pairs contained in both distinguished maximal
families are $[G:N_G(H)]x_H(p,10)$, where
\[
 x_H(5,10)=3024,\qquad x_H(10,10)=2880,
 \qquad [G:N_G(H)]=160005384.
\]
\end{lemma}
\begin{proof}
The complete common generated-subgroup lists for rows 7 and 11
are checked separately for the two types. Every row on either
list lies below row 32. Thus any pair in both families lies in a
row 32 subgroup. Its unique maximal overgroups must be the ones
determined by its order ten entry, so it is their intersection.

An element $s\in H$ of order ten has image the unique involution
in $H/N$: its image cannot be trivial, since $A_6$ has no element
of order ten. The image of $rs$ is central as well, forcing both
entries into $A_6\times C_2$. Write $s=(b,-1)$, with $|b|=5$.
For $p=5$, write $r=(a,1)$ with $|a|=5$; the product is an
involution when $ab$ is an involution or the identity. For $p=10$
the signs are both negative, and only the former possibility
works. The retained literal alternating group calculation gives
$x_{A_6}(5,5)=2880$ and 144 elements of order five. The stated
counts follow. A fresh enumeration of the complete 720-element
product checks them directly.
\end{proof}

\begin{lemma}\label{J}
Within one subgroup $J\cong J_1$, the pairs already contained in
row 7 or row 11 maximal subgroups of $G$ number
\[
 c_J(5,10)=210672,\qquad c_J(10,10)=175560.
\]
\end{lemma}
\begin{proof}
The complete maximal list of $J$ leaves only the types
$A_5\times C_2$, $D_{10}\times D_6$ and $C_{11}:C_{10}$ for a
proper pair of either type. Indeed, the other maximal orders
exclude the pair, except $L_2(11)$, whose element orders exclude
ten. There are 1463 subgroups of type $A_5\times C_2$ in $J$.

The literal $A_5$ calculation gives $x_{A_5}(5,5)=120$; each
such pair generates $A_5$, since a proper subgroup containing
elements of order five lies in a dihedral subgroup of order ten,
where two rotations cannot have product of order two. In
$A_5\times C_2$, there are therefore 120 pairs generating the
whole group for each of our two types. For $(5,10)$ there are
also pairs generating $C_{10}$; for $(10,10)$ there are no such
pairs. All pairs in $D_{10}\times D_6$ of either required type
generate $C_{10}$, or the count is zero. These assertions are
also checked by literal multiplication and generated closures.

In $C_{11}:C_{10}$, every proper pair with an entry of order ten
lies in a cyclic complement of order ten, by its complete maximal
list. Pairs generating the whole order 110 subgroup lie outside
the distinguished families by Lemma~\ref{cover}. Conversely,
$A_5\times C_2$ lies in the centralizer of its central involution,
and every cyclic subgroup of order ten lies in both distinguished
maximal families. Thus the covered pairs in $J$ are precisely
those generating its 1463 subgroups of order 120, together with
those generating a cyclic subgroup of order ten. These are
disjoint by their exact generated subgroups.

The acquired trivial character counts 35112 elements of order
ten in $J$. Hence $J$ has $35112/4$ cyclic subgroups of order ten.
Each supplies four $(5,10)$ pairs and no $(10,10)$ pairs. Therefore
$c_J(5,10)=1463\cdot120+35112$ and
$c_J(10,10)=1463\cdot120$, as claimed.
\end{proof}

\begin{lemma}\label{affine}
Within one row 12 maximal subgroup $M$, the covered pair counts are
\[
 c_M(5,10)=81\cdot384,\qquad c_M(10,10)=81\cdot320.
\]
\end{lemma}
\begin{proof}
Use $M=V\rtimes K$ as in Lemma~\ref{cover}, and let $F$ be the
preimage in $K$ of the rotation subgroup of $K/E\cong D_{10}$.
Then $F=E\rtimes C_5$ has order 160. The roster identifies $K,F$
as rows 184 and 245, and lists 81 conjugate complements $K$ in
$M$. This also follows from conjugacy of complements to the
normal solvable Hall subgroup $V$.

Every element of order five acts without nonzero fixed vectors
on $V$. Its action is nontrivial: otherwise its cyclic subgroup
would be normalized by the order 81 subgroup $V$, contrary to
the global normalizer order 720 of row 574. The polynomial
$X^4+X^3+X^2+X+1$ is irreducible over $\mathbb F_3$, since three
has order four modulo five. A nontrivial action on this
four-dimensional space is consequently irreducible and has no
fixed vectors. The fifth-order part of an order ten element has
this property too. Such an element has a unique fixed point in
the affine action on $V$. The 81 complements are the stabilizers
of these points, so each order ten element lies in a unique one.

The complete common descendant/ancestor calculation shows that
the only possible generated subgroup rows for a pair in $M$
and the distinguished families are $184,245,490,549$. All these
have order prime to three. Each fixes a point in the affine action:
indeed, the average of any orbit is fixed, since its group order
is invertible in $\mathbb F_3$. Thus the actual subgroup inside
the chosen $M$ lies in a complement $K$. In $K/E\cong D_{10}$, an order ten
element has image of order five: its factor five cannot disappear
in the 2-group kernel, and $D_{10}$ has no element of order ten.
An order five element also maps into the rotations. Thus both
entries lie in $F$. Conversely, $F$ centralizes the central
involution of its extraspecial normal subgroup, so is contained
in a row 7 maximal subgroup. We have proved that the covered
pairs are exactly those in the 81 subgroups $F$, with no overlap
because their order ten entries determine their complements.

It remains to count in $F=E\rtimes C_5$. Put $Z=Z(E)=\langle z\rangle$.
The square map defines a nondegenerate quadratic form $q$ on
$E/Z\cong\mathbb F_2^4$. The order five action is nontrivial:
if it centralized $E$, then $E\leq N_G(C_5)$, impossible since
32 does not divide 720. An automorphism of odd order that acts
trivially on $E/Z$ acts trivially on $E$: after fixing the quotient
and $z$, its action multiplies each lift by $z$ or $1$, a group of
exponent two. Hence the action on $E/Z$ is nontrivial and
irreducible, since two has order four modulo five.

A nondegenerate quadratic form in four dimensions over
$\mathbb F_2$ has either five or nine nonzero singular vectors.
For completeness, a symplectic basis writes it as the orthogonal
sum of two binary forms, each with three or one zeros; counting
equal values in the two summands gives ten or six zeros in total.
The fixed-point-free order five action partitions the nonzero
singular vectors into orbits of length five. Their number is
therefore five.

The quotient $F/Z$ is $\mathbb F_2^4:C_5$. Each of its 64 elements
of order five has exactly one lift of order five and one of order
ten. Fix an order five lift $a$. The 16 order five lifts $b$ whose
images in $F/E$ are inverse to that of $a$ make $ab$ run through
one lift of each vector in $E/Z$. Above zero the product is $1$,
since $b=a^{-1}$ is the unique order five lift; above each of the
five nonzero singular vectors it is an involution. Above the
remaining vectors it has order four. Thus $a$ and $zb$ form six
$(5,10)$ pairs, whereas $za$ and $zb$ form five $(10,10)$ pairs.
Multiplication by the 64 choices of $a$ gives 384 and 320.

An independent finite check realizes $E$ by four involutions
whose distinct pairs have commutator $z$, and cycles the five
involutions $g_1,g_2,g_3,g_4,g_1g_2g_3g_4$. It verifies the group
laws, center, square form and all pairs in the resulting group
of order 160. The general count above does not depend on identifying
an unlabelled group by numerical coincidences.
\end{proof}

\begin{theorem}
The ordered automorphism-orbit counts are
\[
 \frac{e_G(5,10)}{2|G|}=365,\qquad
 \frac{e_G(10,10)}{2|G|}=3613.
\]
\end{theorem}
\begin{proof}
The acquired maximal character evaluations have nonzero
contributions only in rows 6, 7, 11 and 12. In the sum over all
maximal subgroups, pairs in the two distinguished families are
counted once in each family and their intersection is corrected
by Lemma~\ref{H}. For $J$ and $M$, remove exactly the covered
contributions from Lemmas~\ref{J} and \ref{affine}. The remaining
pairs have unique maximal overgroups by Lemma~\ref{cover}. Hence
\[
 e_G(p,10)=x_G(p,10)-\sum_T[G:T]x_T(p,10)
 +160005384x_H(p,10)+2624832c_J(p,10)+17778376c_M(p,10).
\]
Here $T$ ranges over the thirteen maximal classes, including zero
contributions. Exact substitution gives the stated integers.
The action of $\operatorname{Aut}(G)$ on generating pairs is free;
its order $2|G|$ is the previously acquired ATLAS premise. The
pairs are ordered and no division by duality is made.
\end{proof}

\paragraph{Attribution and scope.}
The two numerical assertions originate in Connor and Leemans,
\emph{Algorithmic enumeration of regular maps}, Table 1. We retain
their values. The roster, ATLAS structures and character values
are explicit borrowed inputs. The present argument counts the
covered pairs in the affine complement and reuses the earlier
maximal-overgroup and alternating-group proofs. It makes no
priority claim for its elementary group theory, and supplies no
new chirality assertion, Lean certificate or whole-source credit.
\end{document}
